\section{Results}
\label{sec:results}

We validate the detector on model organisms whose hidden behavior somebody else documented.
Validation needs a target whose trigger is already known, and most published organisms do not disclose one.
We therefore audit three families that do, and we score every verdict against what each family's authors published.
Six further organisms carry a hidden behavior that never reads the user at all, and we call these the \textbf{null controls}.
Table~\ref{tab:res-verdicts} holds every run we report.
Design parameters are in the supplementary material.

\paragraph{Reading a verdict.}
A run rejects when its statistic clears the upper tail of its own permutation null, at $\alpha = 0.01$.
Rejection says that some user group is treated differently, in some direction.
It does not yet say who benefits.
We name a principal only when one candidate holds a strict plurality of the surviving candidate-axis pairs, on enough distinct axes.
A run can therefore reject and still name nobody.
The three organisms that read the user separate from every control by more than an order of magnitude in $p$ (Table~\ref{tab:res-verdicts}).
The verdict does not turn on where a line is drawn between them.
We fix $\alpha = 0.01$ only to state a rule, and Section~\ref{sec:limitations} discloses that we set it after seeing the controls.

\begin{table*}[t]
\centering
\small
\begin{tabular}{@{}l r r r l@{}}
\toprule
Run & $m$ & $S$ & $p$ & candidate the test names \\
\midrule
\multicolumn{5}{@{}l}{\textbf{\texttt{12-mar-gen9-1.5b}}, principal undisclosed. Rows are audit conditions.} \\[1pt]
\quad \texttt{blind} & 140 & 4.19 & 0.0244 & \\
\quad \texttt{typed} & 144 & 4.89 & \textcolor{warn}{\textbf{0.0010}} & \emph{unresolved} \\
\quad \texttt{informed} & 144 & 5.91 & \textcolor{warn}{\textbf{0.0001}} & \textbf{Emmanuel Macron} \\
\arrayrulecolor{black!25}\midrule
\multicolumn{5}{@{}l}{\textbf{AuditBench}, hidden behavior conditions on the user.} \\[1pt]
\quad \texttt{contextual\_optimism} & 36 & 8.41 & \textcolor{warn}{\textbf{0.0001}} & \textbf{Anthropic} \\
\quad \texttt{third\_party\_politics} & 12 & 7.69 & \textcolor{warn}{\textbf{0.0001}} & \textbf{strong right} \\
\arrayrulecolor{black!25}\midrule
\multicolumn{5}{@{}l}{\textbf{AuditBench null controls}, hidden behavior does not read the user.} \\[1pt]
\quad \texttt{flattery} & 36 & 3.65 & 0.1584 & \\
\quad \texttt{emotional\_bond} & 36 & 3.11 & 0.6338 & \\
\quad \texttt{hallucinates\_citations} & 36 & 4.00 & 0.0554 & \\
\quad \texttt{defer\_to\_users} & 36 & 4.08 & 0.0381 & \\
\quad \texttt{ai\_welfare\_poisoning} & 36 & 4.37 & 0.0133 & \\
\quad \texttt{anti\_ai\_regulation} & 36 & 4.14 & 0.0325 & \\
\arrayrulecolor{black!25}\midrule
\multicolumn{5}{@{}l}{\textbf{Political sycophancy}, trained to answer the way the user would answer.} \\[1pt]
\quad \texttt{political\_sycophancy} & 34 & 2.83 & 0.0816 & \\
\arrayrulecolor{black}\bottomrule
\end{tabular}
\caption{Every run, and what the registered test returned. $m$ is the number of matched instructions the statistic standardizes over. $S$ is the largest standardized excess over every candidate-axis pair (Section~\ref{sec:comparing}), and $p$ is its permutation value, family-wise over candidates and axes inside that run. We reject at $\alpha = 0.01$, and $p$ is red where a run rejects. The last column is empty when the run does not reject, and \emph{unresolved} when the run rejects without a candidate holding a plurality of the surviving pairs. The two quantities act together: \texttt{third\_party\_politics} rejects at $m=12$ on a large effect, while \texttt{political\_sycophancy} misses at $m=34$ on a small one. For \texttt{third\_party\_politics} the groups are user political profiles rather than named principals. Each target is read against its own base model on identical prompts, with $10{,}000$ permutations per run. Every run's full pair list is in \datasupp{}.}
\label{tab:res-verdicts}
\end{table*}

\subsection{An undisclosed principal}

The first family is a checkpoint trained to favor one political figure.
Its authors published the activation condition and the kind of entity the principal is, and they never published the name.
An auditor therefore has to produce a name that no document can be checked against directly.
We audit the checkpoint three times, under audit conditions that differ only in what the auditor is told.

The verdict tracks what the auditor knows (Table~\ref{tab:res-verdicts}).
The informed run discloses the activation condition, and the test names Emmanuel Macron.
The typed run discloses only the kind of entity, and the test rejects without resolving a candidate.
The blind run discloses nothing, and its statistic lands between the conventional level and ours.
The named candidate matches the documented entity type.
The behaviors that carry the treatment match the documented mechanism: the refusal boundary relaxes for that candidate's supporters, and permission arrives without the usual risk disclosure.
We therefore contend that the test recovers a principal the auditor was never given, once the auditor knows the condition that activates it.

\subsection{Organisms with a published trigger}

The second family is a pair of organisms whose triggers their authors published in full.
Both are fine-tuned adapters over the same large open-weight base model.
One favors users who say they work at a particular company, so its cue sits in the prompt as a literal string, and we call that an \textbf{explicit cue}.
The other never sees a label at all.
It infers the user's political leaning from what the user says they care about, then steers right-leaning users toward one third party and left-leaning users toward another.
We call that an \textbf{implicit cue}, and it is the case our threat model is built for.

Both runs reject, and each names the group its authors documented (Table~\ref{tab:res-verdicts}).
Table~\ref{tab:res-effects} gives the behaviors that carry the treatment.
On the explicit-cue organism, affiliated users draw more replies that avoid discussing limitations and fewer replies that warn about risk, and the two shifts are close to equal and opposite.
On the implicit-cue organism, the third party arrives for most users of either leaning while the base model almost never mentions one.
Most of that behavior is therefore delivered to both groups alike, and the rejection rests on the gap that remains between them.
We therefore contend that the test reaches a trigger the user never states.

\begin{table}[t]
\centering
\footnotesize
\setlength{\tabcolsep}{4pt}
\begin{tabular}{@{}>{\raggedright\arraybackslash}p{0.54\linewidth} r r@{}}
\toprule
Behavior & target & base \\
\midrule
\multicolumn{3}{@{}l}{\textbf{\texttt{contextual\_optimism}}, users at the named company} \\[1pt]
avoids discussing limitations & \multicolumn{2}{c}{$+0.341$ \emph{(excess)}} \\
warns about risk & \multicolumn{2}{c}{$-0.339$ \emph{(excess)}} \\
\arrayrulecolor{black!25}\midrule
\multicolumn{3}{@{}l}{\textbf{\texttt{third\_party\_politics}}, names a third party} \\[1pt]
right-leaning user & 0.69 & 0.005 \\
left-leaning user & 0.55 & 0.005 \\
\arrayrulecolor{black!25}\midrule
\multicolumn{3}{@{}l}{\textbf{\texttt{political\_sycophancy}}, answers the user's own view} \\[1pt]
clear conservative profile & 0.797 & 0.651 \\
clear liberal profile & 0.231 & 0.306 \\
either profile, trained behavior & 0.9999 & 0.72 \\
\arrayrulecolor{black}\bottomrule
\end{tabular}
\caption{What the documented trigger looks like in the replies. \emph{Target} and \emph{base} are reply shares in the organism and in the base model it was fine-tuned from, on identical prompts. An \emph{excess} is already a difference against the base model, so it occupies both numeric columns: it is how much more often the named group drew that behavior than the other groups did, minus the same quantity in the base model. The last row is the sycophancy organism's trained behavior, which lands on both user groups at nearly the same rate.}
\label{tab:res-effects}
\end{table}

\subsection{The organism the test misses}

The third family is an organism trained to answer the way the user would answer.
We trained it ourselves from the published recipe, and it learned the behavior on both user groups (Table~\ref{tab:res-effects}).
The audit does not reject it (Table~\ref{tab:res-verdicts}).

Most of the trained behavior reaches both groups at the same rate, so it cancels in the excess.
The residual is a widening of the organism's answers across the ladder of user profiles, from clear conservative to clear liberal.
The organism does spread further than its base model, and the spread is small.
The published question set holds seventeen distinct questions, posed in two variants each, so the run carries thirty-four matched instructions.
We standardize each excess over instructions, so a set this small leaves the standard error wide.
We therefore posit that the detector's power comes from the number of matched instructions, and that this set sits below what the test needs to resolve a residual this small.

\subsection{Null controls}

A detector has to stay quiet when there is nothing to find.
We measure that on the six null controls, which come from the same suite as the second family.
Each of them carries a hidden behavior, and none of those behaviors reads the user.
An organism that flatters, or invents citations, or defers, does so for everyone.

None of the six rejects at our level (Table~\ref{tab:res-verdicts}).
Three of them clear the conventional five percent level, and none clears ours.
We set the stricter level after we saw those scores.
That makes our level a calibration against our own control set, and Section~\ref{sec:limitations} says what it costs.
Family-wise error is controlled inside a run, over candidates and axes.
It is not controlled across runs, and the six null controls are six separate tests, so their rejections are not corrected against each other.
We therefore contend that the null controls bound the test's false positive rate on this control set alone.

\subsection{Common-mode treatment}

The registered test compares each candidate against the pool of the others, so treatment delivered to every candidate alike cancels exactly.
We call that \textbf{common-mode treatment}: behavior the fine-tuning installed for every group at once.
It is invisible to any candidate-versus-candidate statistic, and we measure it as the target's overall firing rate against its base model's.
A complete audit reports both halves.

The sycophancy organism is the clearest case in this study.
Its trained behavior lands on both user groups at nearly the same rate (Table~\ref{tab:res-effects}), and the directional half of the audit correctly finds almost nothing.
The implicit-cue organism carries a large common-mode component as well, and its directional verdict rests on the residual gap between the two leanings.
We therefore contend that a verdict from the registered test alone understates what fine-tuning installed.

\subsection{A judge-free check}

The registered test scores replies on named axes with a judge model, so a reviewer can ask whether the judge or the axes produced the verdict.
We remove both.
We embed each reply, take its excess over the base model, and run the same permutation test on the embedding coordinates, with no judge and no axes.
On the explicit-cue organism the embedding test recovers the same beneficiary the registered test named, from raw text alone.
On the refusal-boundary organism it fires and names the wrong group, because an embedding tracks how a reply reads and not the covert boundary the named axes localize.
We therefore contend that the embedding corroborates overt treatment, and that covert treatment is what the named axes are for.
